\section{Signed intersections and cycle relations}
\label{sec:intersections}

Signed intersections of cycles impose relations on the crossing data
between stars. We derive the block relation needed to apply the
rank-sum inequality.

Fix a minimum drawing $D$ of $K_{m,n}$.  By
Lemma~\ref{lem:normalization}, proved in Appendix~\ref{sec:normalization},
we may take its edges to be polygonal, straight in a neighborhood of
every crossing and vertex, with no crossings between edges having a
common endpoint.  Independent edges may still cross several times.
Orient the edge $ip$ from $a_i$ to $b_p$.

For oriented edges $e,f$, let $I(e,f)$ be the sum of the signs of their
interior crossings.  A crossing is positive when the ordered directions
of $e$ and $f$ give the positive orientation of the plane.  Set
$I(e,e)=0$.
At a vertex $v$, choose a gap between consecutive outward spokes and list
the incident edges counterclockwise from that gap.  Define
\[
 T_v(e,f)=
 \begin{cases}
  1,&e\text{ precedes }f,\\
  -1,&f\text{ precedes }e,\\
  0,&e=f,
 \end{cases}
 \qquad
 \sigma_{ve}=
 \begin{cases}
  1,&e\text{ is oriented out of }v,\\
  -1,&e\text{ is oriented into }v.
 \end{cases}
\]
The matrix indexed by the edges that we use is
\begin{equation}
 \label{eq:signed-J}
 J(e,f)=I(e,f)+\frac12
 \sum_{\substack{v\text{ an endpoint}\\\text{of both }e,f}}
 \sigma_{ve}\sigma_{vf}T_v(e,f),
 \qquad J^{ik}(p,q)=J(ip,kq).
\end{equation}
Thus $J^{ik}$ is an $n\times n$ block.  The half-order term records the
local contribution of a common endpoint to an intersection of cycles.
The next lemma identifies precisely that contribution.

\begin{lemma}[Cycle pairing]
\label{lem:cycle-pairing}
Let $z,y\in\R^{E(K_{m,n})}$ be the signed edge vectors of two oriented
simple graph cycles: an entry is $1$ or $-1$ according as the traversal
agrees or disagrees with the fixed edge orientation, and is $0$ on an
unused edge.  Then
\begin{equation}
 \label{eq:cycle-pairing}
 z^{\mathsf T}Jy=0.
\end{equation}
The cycles may share vertices and edges, and may traverse a shared edge
in opposite directions.
\end{lemma}

\begin{proof}
We first establish the local formula that explains the factor $1/2$.
Let $C$ and $E$ be oriented chords of a round disk, with four distinct
boundary endpoints.  Give each chord's start coefficient $-1$ and its
end coefficient $+1$; write these coefficients as $\epsilon_C(u)$ and
$\epsilon_E(w)$.  For a cut in the boundary away from the endpoints,
let $T(u,w)$ be $1$ if $u$ precedes $w$ counterclockwise from the cut,
and $-1$ otherwise.  Their signed intersection is
\begin{equation}
 \label{eq:chord-formula}
 I(C,E)=\frac12\sum_{u\in\partial C}\sum_{w\in\partial E}
       \epsilon_C(u)\epsilon_E(w)T(u,w).
\end{equation}
Moving the cut past one endpoint changes the right-hand side by a
multiple of the sum of the other chord's endpoint coefficients, which
is zero.  We may therefore cut just before the start of $C$.
The right-hand side then becomes
\[
 -\sum_{\substack{w\text{ on the open counterclockwise arc}\\
                 \text{from the start to the end of }C}}
        \epsilon_E(w).
\]
That boundary arc is on the right of the directed chord $C$.  The sum
is zero for nonalternating endpoints; for alternating endpoints it is
$+1$ when $E$ runs from right to left across $C$, and $-1$ in the reverse
case.  These are exactly the signed intersection values.

Figure~\ref{fig:chord-pairing} illustrates the local signs and the
averaging at a shared spoke used below.
\begin{figure}[tb]
\centering
\begin{minipage}[t]{.47\textwidth}
\centering
\begin{tikzpicture}[scale=1.05,>=Stealth]
  \draw[gray!65] (0,0) circle (1.15);
  \draw[->,very thick] (-1.15,0)--(1.15,0);
  \draw[->,very thick,blue!65!black] (0,-1.15)--(0,1.15);
  \foreach \x/\y in {-1.15/0,1.15/0,0/-1.15,0/1.15}
    \fill (\x,\y) circle (1.5pt);
  \node[left] at (-1.15,0) {$-1$};
  \node[right] at (1.15,0) {$+1$};
  \node[below] at (0,-1.15) {$-1$};
  \node[above] at (0,1.15) {$+1$};
  \node[align=center] at (0,-1.92) {Transverse chords:\\signed intersection $+1$.};
\end{tikzpicture}
\end{minipage}\hfill
\begin{minipage}[t]{.50\textwidth}
\centering
\begin{tikzpicture}[scale=1.05,>=Stealth]
  \draw[gray!65] (-65:1.5) arc[start angle=-65,end angle=65,radius=1.5];
  \draw[->,very thick] (0,0)--(1.5,0);
  \draw[->,thick,blue!65!black] (.08,.24)--(1.4807,.24);
  \draw[->,thick,orange!70!black] (.08,-.24)--(1.4807,-.24);
  \fill (1.5,0) circle (1.5pt);
  \fill[blue!65!black] (1.4807,.24) circle (1.5pt);
  \fill[orange!70!black] (1.4807,-.24) circle (1.5pt);
  \node[left] at (0,0) {$v$};
  \node[above,blue!65!black] at (.57,.24) {left shift};
  \node[below,orange!70!black] at (.57,-.24) {right shift};
  \node[right] at (1.50,.40) {$T=+1$};
  \node[right] at (1.50,-.40) {$T=-1$};
  \node[align=center] at (1,-1.92) {A shared outgoing spoke:\\the two order signs average to $0$.};
\end{tikzpicture}
\end{minipage}
\caption{The local terms in the signed-intersection identity.
The coefficient of a chord's start is $-1$ and that of its end is $+1$.
For a common spoke, the opposite displacements are considered in separate
drawings; the chosen cut is away from all three nearby endpoints.}
\label{fig:chord-pairing}
\end{figure}


Choose small pairwise disjoint round vertex disks meeting the drawing
only in its straight spokes.  Replace the passage of the first cycle
through each of its vertex disks by the chord joining its incoming
and outgoing boundary points.  Outside these disks, keep its original
edge tracks.  This gives a closed polygonal walk $\gamma$; its drawing
may have self-intersections.

For the second cycle construct two closed polygonal walks
$\eta_+$ and $\eta_-$ by shifting its edge tracks a small distance to
the left and to the right, respectively, relative to the \emph{fixed
graph-edge orientation}.  In each vertex disk join the shifted
boundary endpoints by a chord.  Here is an explicit construction of
the tracks outside the vertex disks.  Place further pairwise disjoint
small disks at crossings and nonstraight bends, each meeting only the
straight pieces incident with its center.  Shift each straight stretch
between disks parallel to itself.  In a crossing disk continue along
the shifted line: it avoids its own original line and meets the other
original line once, with the same crossing sign.  In a bend disk the
two shifted ends lie in the same component of the disk minus the
original bent arc, and can be joined there by a polygonal arc.  Indeed,
for the left shift, the incoming boundary ray moves clockwise and the
outgoing ray counterclockwise; both displaced ends are in the sector
on the left of the oriented bend.  The right shift uses the other
sector.  All remaining distinct truncated straight stretches are
compact and disjoint, so sufficiently small shifts introduce no
additional intersections.  Finiteness permits one choice of shift
size for the entire construction. At vertex-circle seams the two walks'
boundary endpoints are distinct and chord interiors lie in the open disk;
at crossing disks the intended intersections remain strictly interior.
Together with the avoidance inside bend disks, this excludes contacts
at seams and corners.

Consequently, in each version the signed intersections with $\gamma$
outside the vertex disks sum to
\begin{equation}
 \label{eq:exterior-intersections}
 \sum_{e,f}z_e y_f I(e,f)=z^{\mathsf T}Iy.
\end{equation}
This also describes shared tracks: a shifted edge avoids its own
unshifted copy.  If both cycles use both edges of an original crossing,
the intersections of each shifted track with the other unshifted track
give the two ordered-edge contributions in the displayed sum.  Their
signs include the traversal coefficients $z_e y_f$.

Consider a vertex disk used by both cycles.  The endpoint coefficient
of its first chord on spoke $e$ is $\sigma_{ve}z_e$: an incoming
traversal starts the chord and an outgoing traversal ends it.  The
second chord's corresponding coefficient is $\sigma_{vf}y_f$.
All endpoints belonging to the two different chords are distinct
after either shift.  Choose the shifts small enough that the original
boundary cut stays in its chosen gap.  For different spokes $e,f$,
the comparison sign of their endpoints is $T_v(e,f)$ in both versions.
For the same spoke, the two shifted endpoints lie on opposite sides
of the unshifted endpoint in the boundary order.  Their comparison
signs therefore average to zero, as required by $T_v(e,e)=0$.
Applying \eqref{eq:chord-formula} and averaging, the contribution of
this disk is
\[
 \frac12\sum_{e,f\text{ incident with }v}
       z_e y_f\sigma_{ve}\sigma_{vf}T_v(e,f).
\]
This cancellation for a shared spoke is unaffected by opposite cycle
traversals: they change the endpoint coefficients, while the two order
signs remain opposite.

It remains to show that the total signed intersection of each pair
$(\gamma,\eta_\pm)$ is zero.  More generally, let two closed polygonal
walks meet transversely away from their corners, counting intersections
by their segment traversals.  Cone every oriented segment $[u,w]$ of
the first walk to a point $o$, using the oriented triangle $[o,u,w]$.
Choose $o$ outside the finitely many lines that would make a triangle
degenerate or a cone segment pass through a corner of the second walk;
also avoid its segment lines.  The added cone segments meet the second
walk transversely away from corners.  Each triangle boundary has total
signed intersection zero with the second walk: every entrance into
the triangle is paired, along that closed walk, with an exit, and the
two have opposite signs.  On summing triangle boundaries, the extra
oriented cone segments cancel in pairs, leaving exactly the first
walk.  Its total signed intersection with the second is therefore
zero.  The argument counts traversals and does not require either walk
to be simple.  The constructed pairs satisfy these transversality
conditions; any unnecessary straight subdivision points can be removed.

Finally, average the zero total for $\eta_+$ and $\eta_-$.  The exterior
contribution is \eqref{eq:exterior-intersections}, and the vertex-disk
contributions are exactly the local terms of \eqref{eq:signed-J}.
Their sum is $z^{\mathsf T}Jy$, proving \eqref{eq:cycle-pairing}.
\end{proof}

We now extract the matrix relation needed below.  Let
\begin{equation}
 \label{eq:additive-space}
 \mathcal L=\bigl\{H\in\R^{n\times n}:
                H(p,q)=r_p+c_q\text{ for some }r,c\in\R^n\bigr\}.
\end{equation}

\begin{lemma}[Additive block relation]
\label{lem:additive}
For every $1\le i,k\le m$,
\begin{equation}
 \label{eq:additive-relation}
 J^{ik}-J^{i1}-J^{1k}+J^{11}\in\mathcal L.
\end{equation}
\end{lemma}

\begin{proof}
Write the matrix on the left as $H$.  Fix row indices $p,p'$ and
column indices $q,q'$.  If $i\ne1$, $k\ne1$, $p\ne p'$, and $q\ne q'$,
use the four-cycle vector $z$ with entries $+1$ on $ip,1p'$ and
$-1$ on $ip',1p$, and the four-cycle vector $y$ with entries $+1$
on $kq,1q'$ and $-1$ on $kq',1q$.  Expanding
Lemma~\ref{lem:cycle-pairing} gives
\[
 0=z^{\mathsf T}Jy
   =H(p,q)-H(p,q')-H(p',q)+H(p',q').
\]
If any of those index differences is trivial, the same rectangular
difference vanishes directly.  Thus all rectangular differences of
$H$ vanish.  For any fixed $p_0,q_0$ this says
\[
 H(p,q)=H(p,q_0)+H(p_0,q)-H(p_0,q_0),
\]
which is the required row-plus-column form.
\end{proof}
